\documentclass[10pt,a4paper]{amsart}
\usepackage[margin=19mm]{geometry}
\usepackage{amsmath,amssymb,mathtools}
\usepackage[scr=rsfs,cal=euler]{mathalfa}
\usepackage{xfrac}
\usepackage[unicode,hidelinks,bookmarks=false]{hyperref}
\newcommand{\ie}{i.e.\ }
\newcommand{\cf}{cf.\ }
\newcommand{\resp}{respectively\ }
\newcommand{\Subsection}{\subsection}
\providecommand{\nobreakdash}{\nobreak\hbox{-}}
\setlength{\emergencystretch}{2em}
\multlinegap0pt
\usepackage{amssymb}
\usepackage{tikz}
\usepackage{xcolor}
\newtheorem{lemma}{Lemma}
\newtheorem{claim}{Claim}
\title{Signless Laplacian index conditions for trebly chorded cycles in graphs with given order}
\author{Jin Cai, Bo Zhou}
\date{Source: arXiv:2604.15649v1 (CC BY 4.0)}
\begin{document}
\maketitle

\noindent\emph{Source and licence.} Signless Laplacian index conditions for trebly chorded cycles in graphs with given order. Version arXiv:2604.15649v1 (CC BY 4.0), distributed by arXiv under the Creative Commons Attribution 4.0 International licence, http://creativecommons.org/licenses/by/4.0/, as recorded in the arXiv licence metadata for this exact version and shown on the abstract page https://arxiv.org/abs/2604.15649v1. Full licence notice: \texttt{licenses/arxiv-2604.15649-CC-BY-4.0.txt}.\par\smallskip

\noindent\textit{Editorial scope.} Counted unit: the article's claim nok4 (source0.tex lines 1019--1024). The article's complete original proof of it is delivered with it (source0.tex lines 1025--1104, one \texttt{proof} environment of 80 lines). No mathematics is written, repaired or re-derived: the statement and the proof are the article's own lines, in the article's own notation, and no retained line is substituted or re-flowed.  Retained uncounted as the source-local context the proof needs: lines 169--171; lines 178--181.  Nothing was omitted from any retained span, so the package carries no omission record.  No retained line contains a citation, so the wrapper carries no bibliography and no bibliography entry.  No retained line contains a figure environment, an image-inclusion command or drawing code, so no image is delivered and the package declares no asset dependency.  Editorial formatting only, in addition: an AMS \texttt{amsart} preamble that restates the article's own declarations and macros used by the retained lines, namely the environments \texttt{lemma}, \texttt{claim} on separate counters, together with the standard packages those lines need.  The retained environments print in this document as Lemma 1, Claim 1 in order of appearance, whereas the article prints the same items with the numbering of its own full document.  The complete native source of the article as distributed in the arXiv e-print for this exact version is bundled unchanged as \texttt{sources/198628/source0.tex}.
\begin{lemma}\label{addedges}
Let $G$ be a graph and $H$ be a subgraph of $G$. Then $q(H)\le q(G)$. Moreover, if $H$ is a proper subgraph of $G$ and $G$ is connected,  then $q(H)<q(G)$.
\end{lemma}

\begin{lemma} \label{perron}
Let $G$ be a graph with $\{u,v\}\subset V(G)$ and $\emptyset\ne S\subseteq N_G(v)\setminus N_G[u]$. Let
$G'=G-\{vw: w\in S\}+\{uw: w\in S\}$. If $\mathbf{x}$ is the Perron vector of $G$ with $x_u\ge x_v$, then $q(G')>q(G)$.
\end{lemma}

\begin{claim}\label{nok4}
There is at most one component that does not contain $C_4$ in $G[Z^*]$, and any other component
%which contains $C_4$
is $K_4$.
Moreover, if there is one component that does not contain $C_4$ in $G[Z^*]$, then the component is $K_1$, $K_{1,1}$ or $K_{1,s}^+$ for some $s\ge 2$.
\end{claim}

\begin{proof}
We first claim that there is no double star component in $G[Z^*]$. For, suppose that  $S_{b_1,b_2}$
is a  component in  $G[Z^*]$, in which  $v_1$ and $v_2$ are the centers, and $v_{i,1},\dots,v_{i,b_i}$ are the leaves adjacent to $v_i$ for any $i=1,2$.
Assume that $x_{v_1}\ge x_{v_2}$.
Let $G'=G-\{v_2v_{2,j}:1\le j\le b_{2}\}+\{v_1v_{2,j}:1\le j\le b_{2}\}$.
Obviously, $G'$ does not contain a trebly chorded cycle with chords incident to a vertex.
From Lemmas \ref{addedges} and \ref{perron}, we have $q(G)<q(G')$, a contradiction.

Note that if $uv\in E(G)$ and $N[u]\subseteq N[v]$, then since
\begin{align*}
q(G)x_u=d(u)x_u+\sum_{w\in N(u)}x_w\\
q(G)x_v=d(v)x_v+\sum_{w\in N(v)}x_w,
\end{align*}
we have
\[
\left(q(G)-d(u)+1\right)\left(x_v-x_u\right)=\left(d(v)-d(u)\right)x_u+\sum_{w\in N[v]\backslash N[u]}x_w,
\]
so $x_u\le x_v$. If $K_{1,s}$ is a component of $G[Z^*]$, $u$ is the center and $u_{1},\dots,u_{s}$ are the leaves of $K_{1,s}$, then $x_{u}\ge x_{u_{1}}$. Thus, if $K_{1,s}$ is a component of $G[Z^*]$, then the entry of the Perron vector is maximum  at the center among vertices of $K_{1,s}$. Similarly,
%If  $K_{1,s}^+$ is a component of $G[Z^*]$ with center $v$, in which the other vertices are
% $v_{1},\dots,v_{s}$, where $v_{1}$ and $v_{2}$ are adjacent, then $x_{v}\ge x_{v_{i}}$ where $1\le i\le s$. Thus,
 if  $K_{1,s}^+$ is a component of $G[Z^*]$ then the entry of the Perron vector is maximum at the center among the vertices of $K_{1,s}^+$.


%If $K_{1,s}$ is a component of $G[Z^*]$, $u$ is the center and $u_{1},\dots,u_{s}$ are the leaves of $K_{1,s}$, then, by symmetry,
%we have  $x_{u_{i}}=x_{u_{1}}$ for $2\le i\le s$,
%and from $Q(G)\mathbf{x}=q(G)\mathbf{x}$, we have
%\begin{align*}
%q(G)x_{u}&=\left(s+1\right)x_{u}+\sum_{j=1}^{s}x_{u_{j}}+x_{z}=\left(s+1\right)x_{u}+sx_{u_{1}}+x_{z},\\
%q(G)x_{u_{1}}&=2x_{u_{1}}+x_{u}+x_{z},
%\end{align*}
%so $x_{u}\ge x_{u_{1}}$. Thus, if $K_{1,s}$ is a component of $G[Z^*]$, then the entry of the Perron vector is maximum  at the center among vertices of $K_{1,s}$.


We next claim that there exists at most one star component in $G[Z^*]$.
For, suppose that there are two star components in $G[Z^*]$.
Assume that $x_{u_1}$  is maximum among the centers of all star components of $G[Z^*]$.
Let $G''$ be the graph obtained from $G$ by deleting all edges of the star components of $G[Z^*]$ except the one with center $u_1$ and adding edges between $u_1$ and all arising isolated vertices in $Z^*$.
Obviously, $G''$ does not contain a trebly chorded cycle with chords incident to a vertex.
From Lemmas \ref{addedges} and \ref{perron}, we have $q(G)<q(G'')$, a contradiction.

Now we claim that any star component of $G[Z^*]$ must be $K_1, K_{1,1}$. For,
suppose that $K_{1,s}$ is the star component of $G[Z^*]$ with some $s\ge 2$. Let $v_1^*,v_2^*$ be two leaves of $K_{1,s}$.
Obviously, $G+v_1^*v_2^*$ does not contain a trebly chorded cycle with chords incident to a vertex.
From Lemma \ref{addedges}, we have $q(G)<q(G+v_1^*v_2^*)$, a contradiction.


%If  $K_{1,s}^+$ is a component of $G[Z^*]$ with center $v$, in which the other vertices are
% $v_{1},\dots,v_{s}$, where $v_{1}$ and $v_{2}$ are adjacent, then
%by symmetry, we have $x_{v_{1}}=x_{v_{2}}$ and $x_{v_{3}}=x_{v_{j}}$ for $4\le j\le s$,
%and from $Q(G)\mathbf{x}=q(G)\mathbf{x}$, we have
%\begin{align*}
%q(G)x_{v}&=\left(s+1\right)x_{v}+\sum_{j=1}^{s}x_{v_{j}}+x_{z}
%=\left(s+1\right)x_{v}+2x_{v_{1}}+\left(s-2\right)x_{v_{3}}+x_{z},\\
%q(G)x_{v_{1}}&=4x_{v_{1}}+x_{v}+x_{z},~~q(G)x_{v_{3}}=2x_{v_{3}}+x_{v}+x_{z},
%\end{align*}
%so $x_{v}\ge x_{v_{i}}$ where $1\le i\le s$. Thus, if  $K_{1,s}^+$ is a component of $G[Z^*]$ then the entry of the Perron vector is maximum at the center among the vertices of $K_{1,s}^+$.


Finally, we claim that there exists at most one $K_{1,s}^+$ component in $G[Z^*]$ for some $s\ge 2$.
For, suppose that there are two components $K_{1,s_1}^+$ and $K_{1,s_2}^+$ with centers  $u_1$ and $u_2$, respectively.  Assume that $x_{u_1}\ge x_{u_2}$.
Let $G^*$ be the graph obtained from $G$ by deleting all edges of $K_{1,s_2}^+$ and adding edges between $u_1$ and $V(K_{1,s_2}^+)$.
Obviously, $G^*$ does not contain a trebly chorded cycle with chords incident to a vertex.
From Lemmas \ref{addedges} and \ref{perron}, we have $q(G)<q(G^*)$, a contradiction.

Suppose that there are two components, say $H_1$ and $H_2$,  which do not contain $C_4$ in $G[Z^*]$,
According to the above claims, we can assume that $H_1\cong K_{1,t}^+$ and $H_2\cong K_{1,s}$, where $t\ge 2$ and $s=0,1$. Let $z_i$ be the center of $H_i$ with $i=1,2$.  For $\{i,j\}=\{1,2\}$ with $x_{z_i}\ge x_{z_j}$.
Let $G^{**}$ be the graph obtained from $G$ by deleting all edges of $H_j$ and adding edges between $z_i$ and $V(H_j)$.
Obviously, $G^{**}$ does not contain a trebly chorded cycle with chords incident to a vertex.
From Lemmas \ref{addedges} and \ref{perron}, we have $q(G)<q(G^{**})$, a contradiction.
Thus,  there is at most one component that does not contain $C_4$ in $G[Z^*]$. Moreover,
by the above argument, if there is  one component that does not contain $C_4$ in $G[Z^*]$, then the component is $K_1$, $K_{1,1}$ or $K_{1,s}^+$ where $s\ge 2$.

From Lemma \ref{addedges} and our choice of $G$,  any  component of $G[Z^*]$ which contains $C_4$ must be  $K_4$
%
%
%assume that $H\ncong K_4$.
%Let $G^{(6)}$ be a graph obtained by adding all possible edges in $H$.
%Obviously, $G^{(6)}$ does not contain a trebly chorded cycle with chords incident to a vertex.
%From Lemma \ref{addedges}, we have $q(G)<q(G^{(5)})$, a contradiction.
\end{proof}

\end{document}
